
\subsection{Formalisation du problème}

Soit $\mathcal{X} = \{x_1, x_2, \ldots, x_n\}$ l'ensemble des données d'entrée. 
On cherche à minimiser la fonction objectif :

\begin{equation}
  \mathcal{L}(\theta) = \frac{1}{n} \sum_{i=1}^{n} \ell\bigl(f_\theta(x_i),\, y_i\bigr) + \lambda\,\Omega(\theta)
  \label{eq:objectif}
\end{equation}

où $\lambda > 0$ est le paramètre de régularisation et $\Omega(\theta)$ est le terme de régularisation.

\begin{theorem}[Convergence]
  Sous les hypothèses suivantes : (i) ... ; (ii) ..., l'algorithme converge vers un minimum local en $O(1/\sqrt{T})$ itérations.
\end{theorem}
\begin{proof}
  La preuve découle directement de l'application du lemme~\ref{lem:principal}.
\end{proof}

\begin{lemma}
  \label{lem:principal}
  Si $f$ est $L$-lisse et $\mu$-fortement convexe, alors $\|\nabla f(x)\|^2 \geq 2\mu\,(f(x) - f^*)$.
\end{lemma}

\subsection{Architecture proposée}

\lipsum[8]

% Exemple de figure
\begin{figure}[H]
  \centering
  % \includegraphics[width=\linewidth]{figures/architecture.pdf}
  \fbox{\parbox{0.9\linewidth}{\centering\vspace{2cm} [Insérer la figure ici] \vspace{2cm}}}
  \caption{Architecture du modèle proposé. Les blocs colorés représentent...}
  \label{fig:architecture}
\end{figure}

